%% ============================================================================================
%% DRAFT v2, part 1: Introduction, Background, Related Work.  Numbers are macros from numbers.tex.
%% ============================================================================================

\section{Introduction}
\label{sec:intro}

Agentic coding tools such as Claude Code, GitHub Copilot, Cursor, Gemini CLI and Codex run an \emph{agent loop}: a
language model plans, calls tools, observes the results and continues until a goal is met
\citep{Galster2026Harness,he2025lmase}. What the model sees and may do is assembled by the tool's harness from files
that developers commit alongside their code; Galster et al.\ catalogue eight such configuration mechanisms across five
tools \citep{Galster2026Configuring,Galster2026Harness}. One of them, the \emph{subagent}, turns configuration into
organisation. A subagent is a named delegate: a Markdown file whose YAML frontmatter declares when the parent agent
should hand it work (\code{description}), which tools it may call (\code{tools}) and which model runs it, and whose
body becomes its system prompt \citep{AnthropicSubagentsDoc}. The parent delegates a task, the subagent works in its
own context window, and only its result comes back \citep{hua2025contextengineering2,Liu2026DiveClaudeCode}. Each
file therefore combines two kinds of control: a role description that a model interprets, and a tool allowlist that
the runtime enforces.

\paragraph{Permission and authority.}
The allowlist invites a least-privilege reading \citep{saltzer1975protection}: a reviewer that is not given
\code{Write} cannot write. That reading conflates two things that capability security keeps apart. A
\emph{permission} is the right to invoke a tool directly; \emph{authority} is everything the agent can cause,
including through programs it launches \citep[\S8.1]{miller2006robust}. A shell is such a program. An agent granted
\code{Read} and \code{Bash} but not \code{Write} lacks the permission to write files and, wherever its shell commands
run without per-command approval and without a sandbox, keeps the authority to do so. The vendor's documentation
acknowledges the mechanism: its file-editing deny rules do not cover subprocesses that write files indirectly
\citep{AnthropicPermissionsDoc}. How often specifications withhold file tools while keeping
a shell, whether their text needs that shell, and how agents behave under each combination are empirical questions.

\paragraph{The research gap.}
Adoption studies count configuration artifacts across repositories and tools
\citep{Galster2026Harness,Galster2026Dataset,mohsenimofidi2026context,robbes2026agentic}; content studies characterise
what developers write in context files and editor rules \citep{Chatlatanagulchai2026AgentREADMEs,Jiang2026CursorRules};
security studies examine denylists, skills, MCP servers and the security rules written into context files
\citep{Chen2026Denylist,Liu2026SkillsInTheWild,Hasan2026MCP,Yan2026DoNotDeny}. The adoption studies count subagents
as one mechanism among eight, reporting 450 subagents in 131 repositories \citep{Galster2026Harness} and 997 subagent
files in 273 repositories \citep{Galster2026Dataset}, but none analyses what they grant; a static defect scan of
agent configurations does not resolve subagents' tool grants \citep{Kapner2026ScanningHarness}. No study
has measured, at scale, (i) what subagent specifications grant once the runtime has resolved them, (ii) how often a
withholding coexists with a retained capability that subsumes it, (iii) whether restrictions stated in prose agree
with the enforced grant, or (iv) how agents behave, over repeated runs, under each combination. Our contribution is
this measurement and the method that makes it possible, not the mechanism, which follows from what a shell is.

\paragraph{Approach.}
We study \nSpecsEng{} specifications in \nReposEng{} engineered GitHub repositories, reconstructed at a snapshot by git
object identifier. Instead of parsing the \code{tools} field, we use \emph{executable oracles}: by analogy with test
oracles, which decide whether an observed behaviour is acceptable \citep[Def.~2.4]{barr2015oracle}, an executable
oracle for a configuration artifact obtains a property of the artifact by running the program that consumes it and
reading what that program reports, rather than by re-implementing its semantics from documentation. Our oracles start
the coding tool without a credential and read which agents it loaded and which tools each resolves to. We add LLM
annotation of what the specifications say and a controlled experiment of what agents do.

\paragraph{Findings in brief.}
In the average engineered repository \RQthreeIncoherentEng{} of explicit grants that withhold every file-writing tool
keep an unscoped shell. Specifications meant to inspect (a model-annotated label) hold file-writing tools
\RQtwoRoleGapWrite{} less often than those meant to change code, but the gap shrinks to \RQtwoRoleGapAuthority{} once
the shell is counted. That a shell permits a requested change when commands are approved automatically is expected,
and our mechanism check confirms it (\ExpCtwoRate{} under the most common inspection grant and a neutral prompt).
Three results are less obvious. A prompt restriction on the same grant, the realistic inspection configuration, holds
for Sonnet and Opus but not reliably for Haiku. In a small report-only extension, no agent changed a file it was not
asked to change. And the text of most write-withholding specifications that keep a shell gives no reason to change
anything by command. We also observed that Haiku without a shell sometimes reports a change it did not make.

\paragraph{Research questions.}
\begin{description}\itemsep1pt
  \item[RQ1 (Delegation structure)] How many subagents do engineered repositories define, for which roles, and how
        much of the population is reused? (\S\ref{sec:res-rq1})
  \item[RQ2 (Capability grants)] How much capability do specifications grant, explicitly or by inheritance, and does
        the effective grant track the delegate's intended mode of work? (\S\ref{sec:res-rq2})
  \item[RQ3 (Coherence of restrictions)] When a specification withholds file-writing tools, does it retain a shell
        that subsumes them, and does its text need that shell? (\S\ref{sec:res-rq3})
  \item[RQ4 (Prose versus enforcement)] Which constraints do developers write in natural language, and does the
        enforced grant agree with them? (\S\ref{sec:res-rq4})
  \item[RQ5 (Behaviour under repeated execution)] How do tool grants and prompt restrictions constrain agents across
        tasks and models, and do agents change files they were not asked to change? (\S\ref{sec:res-rq5})
  \item[RQ6 (Silent defects and drift)] Which specifications does the runtime silently discard, mis-resolve or shadow,
        and do unresolvable tool names reflect runtime evolution or stale sources? (\S\ref{sec:res-rq6})
  \item[RQ7 (Governance context)] Are agent grants associated with the repository's other access controls?
        (\S\ref{sec:res-rq7})
\end{description}
Table~\ref{tab:rqmap} maps each question to the data and the methods that answer it.

\begin{table}[t]
\caption{Research questions, the data and methods that answer them (Section~\ref{sec:design}), and where the results are
reported (Section~\ref{sec:results}).}
\label{tab:rqmap}
\footnotesize
\begin{tabular}{@{}l>{\raggedright\arraybackslash}p{3.4cm}>{\raggedright\arraybackslash}p{5.5cm}l@{}}
\toprule
\textbf{RQ} & \textbf{Data} & \textbf{Method} & \textbf{Results} \\
\midrule
RQ1 & specifications and their roots & census oracle (\S\ref{sec:oracles}); role and mode annotation, codebook R (\S\ref{sec:annotation}); content hashes & \S\ref{sec:res-rq1} \\
RQ2 & \code{tools} and \code{disallowedTools} values & capability model (\S\ref{sec:capmodel}); resolver oracle (\S\ref{sec:oracles}); mode labels; paired within-repository contrasts (\S\ref{sec:stats}) & \S\ref{sec:res-rq2} \\
RQ3 & explicit grants; committed settings & dominance (\S\ref{sec:capmodel}); shell-use annotation, codebook B (\S\ref{sec:nlp}); scoped-shell probe (\S\ref{sec:experiment}) & \S\ref{sec:res-rq3} \\
RQ4 & description and body text & directive-sentence lexicon and codebook C; whole-text restriction, codebook S (\S\ref{sec:nlp}); resolver oracle & \S\ref{sec:res-rq4} \\
RQ5 & controlled fixtures & factorial execution experiment, report-only extension, codebook K (\S\ref{sec:experiment}) & \S\ref{sec:res-rq5} \\
RQ6 & names, roots, settings; release history & census, resolver and settings oracles; semantics probes; tool-pool history and drift dating (\S\ref{sec:oracles}) & \S\ref{sec:res-rq6} \\
RQ7 & repository settings and governance files & governance baseline (\S\ref{sec:governance}); differences of repository means (\S\ref{sec:stats}) & \S\ref{sec:res-rq7} \\
\bottomrule
\end{tabular}
\end{table}

\paragraph{Contributions.}
(1) A snapshot-exact corpus of \nSpecsAll{} subagent specifications from \nReposAll{} repositories, every recovered
file verified against its hash, curated to \nSpecsEng{} specifications in \nReposEng{} engineered repositories.
(2) Executable oracles that make the coding tool report what it loads and grants.
(3) Model annotations, with published codebooks, of every specification's role and intended mode, of the write
restriction a specification's text states, and of what the text asks a retained shell to do. (4) The measurement of \emph{dominated withholding}: explicit grants that omit every
file-writing tool but keep an unscoped shell, which leaves the withholding without effect wherever shell commands run
without approval and without a sandbox. (5) A controlled experiment of
\nExpRuns{} scored delegations, and an extension of \nOverRuns{} report-only delegations, separating what a grant
permits from what agents do. (6) A candidate security smell for agent configuration, guidance for frameworks and a
linter.

%% --------------------------------------------------------------------------------------------
\section{Background}
\label{sec:background}

\subsection{Multi-agent architectures and agent profiles}
\label{sec:bg-mas}
An LLM-based \emph{agent} directs its own process and tool use, whereas a \emph{workflow} orchestrates models and tools
along predefined code paths \citep{anthropic2024buildingagents}. He et al.\ describe an LLM-based multi-agent system as
an orchestration platform, which manages coordination, communication and planning, together with the agents it
coordinates, and list the roles that recur in code generation, among them an orchestrator that decomposes goals and
delegates subtasks \citep[\S2.3, \S3.2]{he2025lmase}. Each agent is steered by a \emph{profile} written into its prompt;
profiles are handcrafted, generated by a model or aligned with data \citep[\S2.1.1]{wang2024survey}, and how agents are
profiled is one of the axes along which multi-agent systems are surveyed \citep{guo2024lmmas}. The role-specialised
frameworks handcraft them. MetaGPT specifies each role by name, profile, goal and constraints and gives each role its
own tools \citep[\S3.1]{hong2024metagpt}; ChatDev assigns every role a system prompt that states its objective,
accessible tools, communication protocol, termination conditions and constraints on undesirable behaviour
\citep[\S3.1]{qian-etal-2024-chatdev}; AutoGen configures conversable agents with model, tool and human capabilities
\citep[\S2.1]{wu2024autogen}. In production the dominant pattern is orchestrator--worker: a lead agent spawns
subagents, each with its own context window, and synthesises their results, and each subagent needs an objective, an
output format, guidance on tools and sources, and task boundaries \citep{anthropic2025multiagentresearch}. How agents
are specified matters: system-design issues are one of the three categories of a taxonomy of multi-agent failures
\citep{cemri2025mast}. Hassan et al.\ expect agentic software engineering to rest on versioned, human-authored
artefacts that brief and constrain agents, and read files such as \code{CLAUDE.md} as an early grassroots form of them
\citep[\S5.3]{hassan2025sase}. A subagent specification is the practitioner's version of a handcrafted agent profile,
committed to version control next to the code it works on.

\subsection{Context engineering and context isolation}
\label{sec:bg-context}
Context engineering is the design of everything a model sees at inference time. Anthropic describes it as curating
and maintaining the optimal set of tokens \citep{anthropic2025contextengineering}; Mei et al.\ formalise the context as
assembled from instructions, knowledge, tool definitions, memory, state and the query, and context engineering as
optimising that assembly under a length constraint \citep[\S3.1]{mei2025contextengineering}; Hua et al.\ define it as the
systematic design of how context is collected, stored, managed and used \citep[\S2.1]{hua2025contextengineering2}.
Context is a finite resource. Accuracy falls when relevant information sits in the middle of a long input
\citep{liu2024lost}, most models degrade well before their advertised context length \citep{hsieh2024ruler}, the
degradation is steeper when the question shares little vocabulary with the relevant passage
\citep{modarressi2025nolima}, and performance grows less reliable with input length even on simple, controlled tasks
\citep{hong2025contextrot}. Delegation therefore doubles as \emph{context isolation}: a subagent works in a fresh
context window with its own system prompt and tool set, and only a condensed result returns to the parent
\citep{anthropic2025contextengineering}\citep[\S5.3.2]{hua2025contextengineering2}. Hua et al.\ name Claude Code's
subagents, with their restricted tool permissions, as an instance and recommend giving each unit only the permissions
it needs \citep[\S5.3.2]{hua2025contextengineering2}, and running reusable skills as subagents outperforms loading them
into the main context when their inputs and outputs are clearly specified \citep{piriyakulkij2026subagents}.

A subagent specification is therefore a context-engineering artefact, and it acts along two channels. The
\emph{interpretive channel} (context files, rules, a subagent's description and body) is text loaded into the model's
context, which the model may or may not follow \citep{AnthropicMemoryDoc}. The \emph{enforcement channel} (tool
allowlists, permission rules, hooks, sandboxes) is applied by the harness regardless of what the model intends
\citep{AnthropicPermissionsDoc}.

\subsection{Configuration mechanisms across agentic coding tools}
\label{sec:bg-mechanisms}
Table~\ref{tab:mechanisms} summarises the mechanisms that frame our study. Context files dominate adoption, appearing
in 90.6\% of configured repositories, while skills and subagents are adopted far more shallowly
\citep[\S6.1--6.3]{Galster2026Harness}. The tools expose the same mechanisms but enforce different parts of them.
\emph{Claude Code} loads \code{CLAUDE.md} and \code{.claude/rules} into the context at session start and treats them as
context, not as enforced configuration \citep{AnthropicMemoryDoc}; permission rules in committed settings are evaluated
deny, then ask, then allow, and are enforced by the tool rather than the model \citep{AnthropicPermissionsDoc,
AnthropicSettingsDoc}; hooks bound to lifecycle events can block a tool call, including hooks declared in a subagent's
frontmatter \citep{AnthropicHooksDoc}. \emph{Cursor} applies project rules (\code{.cursor/rules}, with
\code{description}, \code{globs} and \code{alwaysApply}) by including them at the start of the model context, always,
when the agent judges them relevant, for matching files or on request \citep{CursorRulesDoc}; its subagents have their
own context windows, load from \code{.cursor/agents} and, for compatibility, from \code{.claude/agents} and
\code{.codex/agents}, and offer a \code{readonly} flag \citep{CursorSubagentsDoc}. \emph{GitHub Copilot} adds
repository-wide and path-scoped instruction files and \code{AGENTS.md} to its requests
\citep{GitHubCopilotInstructionsDoc}; its custom agents are Markdown profiles in \code{.github/agents} whose \code{tools}
field defaults to all tools and ignores names it does not recognise \citep{GitHubCreateCustomAgentsDoc,
GitHubCustomAgentsDoc}. \emph{Codex} concatenates the \code{AGENTS.md} files from the project root down to the working
directory \citep{OpenAICodexAgentsMdDoc} and defines custom agents as TOML files in \code{.codex/agents}, which may set a
sandbox mode and otherwise inherit the session's sandbox policy \citep{OpenAICodexSubagentsDoc}; \code{AGENTS.md} is an
open format that several of these tools read \citep{AgentsMdSpec}. Subagent formats have thus converged on a named,
described profile whose body is a prompt, while differing in the restriction they enforce: a tool allowlist (Claude
Code, Copilot), a read-only flag (Cursor) or a sandbox mode (Codex). Because Cursor also reads
\code{.claude/agents}, a file in our corpus may be consumed by more than one tool.

\begin{table}[t]
\caption{Configuration mechanisms relevant to subagents, by the channel through which they alter agent behaviour.
Sources: vendor documentation accessed September 2026 and \citet{Galster2026Harness}.}
\label{tab:mechanisms}
\footnotesize
\begin{tabular}{@{}>{\raggedright\arraybackslash}p{2.1cm}>{\raggedright\arraybackslash}p{2.0cm}>{\raggedright\arraybackslash}p{6.9cm}@{}}
\toprule
\textbf{Mechanism} & \textbf{Channel} & \textbf{How it alters behaviour at run time} \\
\midrule
Context files & interpretive & loaded into the context at session start; \code{AGENTS.md} read by Codex, Copilot, Cursor; \code{CLAUDE.md} by Claude Code \\
Rules & interpretive & included at the start of the context when their globs or description match (Cursor, Claude Code \code{.claude/rules}) \\
Skills, commands & interpretive & instructions loaded on invocation or when the model selects them \\
\textbf{Subagents} & \textbf{both} & body and description are interpreted; the tool allowlist, model and permission fields are applied by the harness \\
Settings & enforcement & permission rules evaluated deny, then ask, then allow, outside the model \\
Hooks & enforcement & scripts at lifecycle events that can block a tool call \\
MCP servers & enforcement & external tools added to the pool when the server is configured \\
\bottomrule
\end{tabular}
\end{table}

\subsection{The Claude Code subagent artifact}
\label{sec:bg-artifact}
A subagent is a Markdown file under \code{.claude/agents/} of a project or of the user's home directory whose
frontmatter must declare \code{name} and \code{description} \citep{AnthropicSubagentsDoc}. Table~\ref{tab:schema}
lists the fields and their documented semantics; what the runtime does with them on the build we study is established
by our probes and reported with the results (\S\ref{sec:res-rq6}, Table~\ref{tab:semantics}). Each subagent runs in its
own context window with its own tool access, and only its result returns to the parent. One documented property
carries the study: \textbf{omission is inheritance}. A specification without \code{tools} receives every tool
available to subagents, so omitting the field is the most permissive choice \citep{AnthropicSubagentsDoc}.

\begin{table}[t]
\caption{Subagent frontmatter on the pinned build (Claude Code v2.1.233). ``Observed'' entries were established by
the controlled probes of Section~\ref{sec:oracles}; ``documented'' entries rest on vendor documentation only.}
\label{tab:schema}
\footnotesize
\begin{tabular}{@{}>{\raggedright\arraybackslash}p{2.2cm}>{\raggedright\arraybackslash}p{2.6cm}>{\raggedright\arraybackslash}p{6.2cm}@{}}
\toprule
\textbf{Field} & \textbf{Meaning} & \textbf{Semantics (source)} \\
\midrule
\code{name} & delegation identifier & required; with a non-string value the file is not loaded, without a diagnostic (observed; an absent name is documented to have the same effect). With two files declaring the same name under one root exactly one loads, and which one is not predictable from the file names; when the session starts in a sub-project, its file wins over the parent project's (observed) \\
\code{description} & when the parent delegates & required; absent and the file is not loaded, without a diagnostic (observed) \\
\code{tools} & tool allowlist & \textbf{omitted, \code{"*"}: the full pool} (observed). Names resolve with aliases (\code{Agent}$\to$\code{Task}, \code{KillShell}$\to$\code{TaskStop}, \code{BashOutput}$\to$\code{TaskOutput}); unresolved names are dropped without a diagnostic; \code{Bash(\ldots)} resolves to \code{Bash}; \textbf{\code{Grep} and \code{Glob} are dropped whenever \code{Bash} is granted}; case-sensitive (\code{read} resolves to nothing) (observed) \\
\code{disallowedTools} & subtractive filter & removed from the inherited or listed set (observed) \\
\code{model} & model tier & alias or model identifier (documented) \\
\code{permissionMode} & autonomy level & ignored when the parent runs in \code{bypassPermissions}, \code{acceptEdits} or \code{auto} mode (documented; \citealp[\S8]{Liu2026DiveClaudeCode}) \\
\code{maxTurns}, \code{effort} & budget & enforced (documented) \\
\code{skills}, \code{mcpServers}, \code{hooks} & capability attachment & ignored for plugin subagents (documented) \\
\code{background}, \code{isolation}, \code{memory} & execution policy & documented \\
\code{color} & user-interface colour & no runtime effect (documented) \\
file location & scope & any depth of nesting under \code{.claude/agents/} loads; a sub-project's \code{.claude/agents} is not loaded from the repository root (observed) \\
\bottomrule
\end{tabular}
\end{table}


\subsection{Protection, least privilege and capability lists}
\label{sec:bg-protection}
Saltzer and Schroeder define least privilege as running every program with the minimum privileges its job requires,
and fail-safe defaults as basing access on permission rather than exclusion \citep[\S I-A-3]{saltzer1975protection}.
Lampson's access matrix records the rights each domain holds over each object, and a capability list, such as a
subagent's allowlist, is one row of it \citep{lampson1974protection}. Systems whose subjects exercise authority without
selecting it hold \emph{ambient authority}, which capability systems remove
\citep{miller2003capability,watson2010capsicum}; a shell carries the process's ambient file-system authority. About a third of 940 Android applications were over-privileged because developers misunderstood what permissions
imply \citep{felt2011android}, and the permission specification had to be extracted from the platform because its
documentation was incomplete \citep{au2012pscout}; for the same reason our labels come from executing the tool.

%% --------------------------------------------------------------------------------------------
\section{Related work}
\label{sec:related}

\paragraph{Empirical studies of agent configuration.}
\label{sec:rw-config}
Galster et al.\ study the adoption of eight configuration mechanisms in 2,853 engineered repositories and release a
dataset of 4,738 configured repositories \citep{Galster2026Configuring,Galster2026Harness,Galster2026Dataset}; they
count 450 subagents in 131 repositories \citep[\S6.3]{Galster2026Harness}, and 997 in 273 repositories in the larger
dataset \citep{Galster2026Dataset}, but do not analyse their content beyond noting that none use the persistent-memory field. Mohsenimofidi et al.\ find context files in 5\%
of 10,000 popular, mature repositories and describe instruction styles ranging from descriptive and prescriptive to
prohibitive \citep{mohsenimofidi2026context}. Context files rarely mention security \citep{Chatlatanagulchai2026AgentREADMEs}, and Cursor rules show substantial
copying across repositories \citep{Jiang2026CursorRules}. Skills are copied at scale without a registry
\citep{Destefanis2026GitSkills}, and a scanner flagged at least one vulnerability pattern in about a quarter (26.1\%) of
31,132 marketplace skills \citep{Liu2026SkillsInTheWild}. Closest to our questions, Kapner et al.\ scan 3,171
repositories for supply-chain defects in agent configurations with rules decided from the files' bytes
\citep{Kapner2026ScanningHarness}. The subagent defect they confirm is a file without a description, which the tool skips;
their closest security class is a scoped-looking pre-approval such as \code{Bash(python:*)} that permits arbitrary
execution (3.1\% of setups). They do not resolve tool grants or relate them to a delegate's purpose or text. Yan finds
that only 4--16\% of security rules written in \code{CLAUDE.md} files have a matching built-in control
\citep{Yan2026DoNotDeny}. Relative to these studies, we obtain the effective grant by running its consumer, measure how
often a withholding is subsumed by a retained shell, compare the grant with the specification's intended mode and
prose, and test behaviour by repeated execution.

\paragraph{Context engineering and prompt artefacts in software repositories.}
\label{sec:rw-context}
Mining studies treat prompts and agent instructions as software artefacts. Prompts in GitHub repositories evolve mostly
through additions and modifications during feature work, and only a fifth of prompt changes are documented in commit
messages \citep{tafreshipour2025prompting}; they co-evolve with code without static checking, which motivated a prompt
linter \citep{pister2024promptset}; and open-source prompts show inconsistent formatting and substantial duplication
within and across repositories \citep{li2026promptmanagement}. Developers form mental models of a model's behaviour
under a prompt rather than of code, and struggle to make them reliable \citep{liang2025promptsareprograms}. For agentic
coding, \code{CLAUDE.md} manifests are dominated by operational commands, implementation notes and architecture
\citep{Chatlatanagulchai2025Manifests}, and agent context files evolve like configuration code through frequent small
additions \citep{Chatlatanagulchai2026AgentREADMEs}. Studies of agent activity on GitHub mine the pull requests agents
open \citep{Li2026AIDev}. On the design side, contexts are curated as evolving playbooks \citep{zhang2026ace}. These
studies analyse what the interpretive channel says; none resolves what a subagent is permitted to do.

\paragraph{Agent permissions and the limits of prompt constraints.}
\label{sec:rw-security}
Chen and Lin show that, depending on the file-system operation, 69.0--98.6\% of real command denylists for terminal
coding agents that target that operation overlook at least one working bypass \citep{Chen2026Denylist}; they study
denylists only, not allowlists or subagent tool grants. On benign tasks built to invite out-of-scope actions, removing the
authorised-scope declaration from the prompt raises Claude Code's rate of over-eager actions from 0.0\% to 17.1\% on
paired scenarios, and the framework's permission gating dominates effect size \citep{qu2026overeager}. Tool filtering fails against prompt injection whenever the permitted tools suffice for the attack
\citep{debenedetti2024agentdojo}, which motivates enforcing policy outside the model
\citep{debenedetti2025defeating,shi2025progent} and replacing shells with narrow tools \citep{beurerkellner2025design}.
Models trained to respect an instruction hierarchy obey it only partially under conflict and strong attacks
\citep{wallace2024hierarchy,geng2026controlillusion}.

\paragraph{Configuration errors and security smells.}
\label{sec:rw-smells}
Infrastructure-as-code research names recurring insecure configuration patterns as security
smells, each with a stated rationale and a CWE mapping \citep{rahman2019seven,rahman2021security}, and a Kubernetes
study observes that a privileged security context makes every other access control in that security context obsolete
\citep[\S2.3]{rahman2023security}. In CI, almost all workflows ran with over-privileged tokens because the permissions
field was rarely set \citep{koishybayev2022characterizing}. The pattern we study is the agent-configuration analogue: a
coarse grant (an unscoped shell) that can void a fine-grained withholding in the same policy.

\paragraph{Mining and annotation methodology.}
\label{sec:rw-method}
Most GitHub repositories are personal, inactive or not software development
\citep{kalliamvakou2016promises}; engineered projects have been separated by dimension scores
\citep{munaiah2017curating} and by LLM classification of READMEs \citep{Galster2026Dataset}. We follow the ACM SIGSOFT
Empirical Standards \citep{ralph2021empirical} and the guidelines for studies that use LLMs, which require exact model
versions, published prompts, repeated runs and validation against humans \citep{baltes2026guidelines,ahmed2025llms}.
